% Auto-generated by R32_sensitivity_tuning.ipynb -- check the wording before pasting.

% Numbers are percentage points of ROC-AUC; gap = classical minus family (positive = classical better).

The configuration audit showed that, with the library default \texttt{drop\_last=True} and a batch size of 256, TabNet received no gradient update in folds whose training split had fewer than 256 rows (5 datasets in at least one mode). Re-running TabNet with \texttt{drop\_last=False} changed its mean Mode-4 ROC-AUC from 65.44\% to 67.22\% and the deep tabular gap from 5.75 to 5.15 points (change -0.59, 95\% CI $[-1.98, +0.62]$).

In the sensitivity ladder, the deep tabular gap was 4.85 points with the default configuration (95\% CI $[3.23, 6.72]$). The arm that narrowed it most was \texttt{lr\_plateau} (shift -2.12 points, CI $[-3.33, -1.08]$); 1 of the 8 arms changed the gap significantly after Holm correction, and the gap remained positive with a confidence interval excluding zero under every arm. Even when the best arm was chosen per model and dataset on the test folds, an optimistic bound, the gap was 1.88 points (CI $[0.90, 3.04]$).

In the sensitivity ladder, the Transformer gap was 2.92 points with the default configuration (95\% CI $[1.43, 4.54]$). The arm that narrowed it most was \texttt{batch\_64} (shift -0.91 points, CI $[-1.99, -0.08]$); no arm changed the gap significantly after Holm correction, and the gap remained positive with a confidence interval excluding zero under every arm. Even when the best arm was chosen per model and dataset on the test folds, an optimistic bound, the gap was 0.63 points (CI $[-0.41, 1.80]$).

For FT-Transformer without oversampling, the feature-selection effect was -2.76 points (CI $[-5.57, -0.46]$) with the default configuration and -0.56 points (CI $[-1.58, +0.48]$) with the compact small-data configuration.